\begin{lemma}[Exact prefix repackaging of a good hierarchy]
Let $G$ be a finite loopless multigraph with distinguished uncolored edge
$e_0$, let $[k]$ be the color set, and let $\varphi_{\rm in}$,
$\varphi_{\rm gen}$, and $\varphi$ be partial edge colorings.  Let
$\mathcal E$ be a positive-rung \noderef{ExtendedTashkinovTree} on the
exact ordered tree sequence $T$ with input coloring $\varphi_{\rm in}$ and
generating coloring $\varphi_{\rm gen}$.  Suppose that $\mathcal E$ has
\noderef{MaximumProperty}.  Let $B$ be a \noderef{HierarchyBase} for
$\mathcal E$ under $\varphi$, let $H$ be a
\noderef{TashkinovHierarchy} with divider count $q$, and let $\Gamma$ be a
\noderef{GoodHierarchy} for $H$.

For every integer $r$ with $1\leq r\leq q$, put $T_r=H.\mathit{stage}(r)$.
There exist an ETT $\mathcal E_r$ on the exact ordered tree sequence $T_r$,
a hierarchy base $B_r$, a hierarchy $H_r$, and a good hierarchy
$\Gamma_r$ such that all of the following hold.
\begin{enumerate}
\item $\mathcal E_r$ has the same rung number and the same entire certified
Tashkinov series as $\mathcal E$, including the exact ordered tree sequences
and the individual identities of every connecting edge, and
$\mathcal E_r$ has maximum property.
\item The core, closure, and numerical support position of $B_r$ are exactly
those of $B$.
\item The divider count of $H_r$ is $r-1$.  For every $0\leq j\leq r$,
$H_r.\mathit{stage}(j)=H.\mathit{stage}(j)$ as ordered tree sequences, so
all displayed vertex and individual edge identities agree.  For every
$0\leq i<r-1$, the numerical divider position in $H_r$ is the corresponding
divider position in $H$.
\item For every $0\leq j<r$, the stored-color sets agree exactly,
$D^{H_r}_j=D^H_j$, and for every color $\eta$ the total Gamma data agree,
$\Gamma^j_{r,\eta}=\Gamma^j_\eta$.  In particular, this equality holds for
every paper-relevant $\eta\in D^H_j$.
\item The old last internal closure field is retained separately:
\[
 T_r\text{ is }
 \left(\bigcup_{\eta\in D^H_r}\Gamma^{r-1}_\eta\right)^{-}
 \text{-closed under }\varphi.
\]
\end{enumerate}
All equalities of tree sequences above are equalities of the ordered
edge-valued sequences, not merely equalities of their vertex sets.
\end{lemma}
\begin{proof}
Write the stages of $H$ as $T_0,T_1,\ldots,T_q,T$.  Repeatedly composing
the ordered initial-segment maps in \noderef{TashkinovHierarchy} shows that
$T_a$ is an ordered initial segment of $T_b$ whenever
$0\leq a\leq b\leq q+1$, with $T_{q+1}=T$.  These maps preserve the actual
vertex entries and edge entries at every retained index, by
\noderef{TreeSequenceInitialSegment}.

We first construct $\mathcal E_r$.  Its rung number and certified series
are literally those of $\mathcal E$, and its displayed terminal tree is
$T_r$.  The last ladder tree is $T_0$ by the base-stage field of $H$, so it
is an initial segment of $T_r$.  It is strictly shorter: the strict
containment $T_0\subset T_1$ is followed by $T_1\subseteq T_r$, using
$r\geq1$.  Also $T_r\subseteq T$, and composing this inclusion with the
old upper inclusion of $T$ into the final certified closure gives the
required upper inclusion for $T_r$.  Thus every field of
\noderef{ExtendedTashkinovTree} is satisfied, with the same certified
series and hence with the same exact connecting-edge identities.  The
definition \noderef{MaximumProperty} refers only to that rung and certified
series; therefore the old maximum-property witness is also a witness for
$\mathcal E_r$.

Because the ladder series did not change, the core, closure, support
position, terminal missing-color equation, PE data, and non-PE data of $B$
have exactly the types required for a \noderef{HierarchyBase} of
$\mathcal E_r$.  Copy them to obtain $B_r$.

Define $H_r$ to have divider count $r-1$ and stage list
$T_0,T_1,\ldots,T_r$.  Its base-stage, terminal-stage, nestedness,
strictness, and base-inside-first-stage fields are the corresponding fields
of $H$.  For $0\leq i<r-1$, define its $i$-th divider position to be
$|T_{i+1}|$.  The old divider-stage identity says
\[
 T_{i+1}=\noderef{TreeSequencePrefix}(T,d_i),
\]
where $d_i$ is the old numerical divider position; taking lengths gives
$d_i=|T_{i+1}|$.  Since $T_{i+1}$ is an initial segment of $T_r$, it is
also the exact prefix of $T_r$ of that length.  Hence the new divider-stage
identity holds, the numerical positions are unchanged, and positivity and
strict increase follow from positivity and strict nesting of the retained
stages.  This constructs the claimed \noderef{TashkinovHierarchy} $H_r$
and proves all equalities in item 3 at the level of the complete ordered
tree sequences.

We next record the only point not obtained by copying record fields.  If an
ordered tree sequence $A$ is an initial segment of $C$, a color $\eta$, and
an index $s<|A|$ satisfy
\[
 s<\noderef{HierarchyFirstMissingPosition}(\varphi,A,\eta),
\]
then $s$ is also smaller than the first missing position in $C$.  Indeed,
if the candidate set in $C$ is empty, its totalized first position satisfies
$s<|C|$.  Otherwise, if its least candidate had index at most $s$, that
index would lie in $A$ and the initial-segment vertex identity would make
it a missing-color candidate in $A$.  If the candidate set in $A$ is
nonempty, its minimum would then be at most $s$, a contradiction; if it is
empty, the exhibited candidate is itself a contradiction.

Apply this observation to $A=T_r$ and $C=T$.  At a retained Gamma index
$j<r$, the starred stage and the next stage used by
\noderef{HierarchySegmentColors} are literally the same in $H_r$ and $H$.
Every edge index in the new segment lies before $T_{j+1}\subseteq T_r$.
The observation therefore shows that every color in the new segment-color
set also lies in the old segment-color set.  Notice that only this inclusion
is needed; no recoloring or generic equality of segment-color sets is being
asserted.

Now define $\Gamma_r^j(\eta)=\Gamma^j(\eta)$ for $0\leq j<r$.  Unfolding
\noderef{HierarchyStoredColorsAt}, the certified stored set is unchanged
and the starred stage is unchanged, so $D^{H_r}_j=D^H_j$.  The stage
missing-color sets defined by \noderef{HierarchyStageMissingColors} also
agree.  We verify the five fields of
\noderef{GoodHierarchy}.  Pair cardinality and pairwise disjointness are
the old fields after the stored-set equality.  For avoidance, the old pair
lies in the common stage-missing set and avoids the old segment-color set;
the inclusion proved in the preceding paragraph implies that it avoids the
new segment-color set as well.  For every $i<r-1$, the two stages, the
starred stage, their vertex-set difference, both stored sets, and both
Gamma families occurring in unique refresh are unchanged, so the old
existence-and-uniqueness witness is exactly the new one.  The parallel-base
field uses only the unchanged series, base, zeroth stored set, and zeroth
Gamma family, and hence copies verbatim.  Finally, for $i<r-1$ the stage,
stored set, and preceding Gamma family in the internal minus-closed field
are unchanged, so that field also restricts verbatim.  Thus $\Gamma_r$ is a
good hierarchy and item 4 holds.

It remains to prove item 5.  Since $r\geq1$, the index $r-1$ is a valid old
internal index.  The `minusClosed' field of the original good hierarchy at
that index states exactly that $T_r$ is
\[
 \left(\bigcup_{\eta\in D^H_r}\Gamma^{r-1}_\eta\right)^{-}
\text{-closed}.
\]
This is precisely the predicate \noderef{TreeColorSetMinusClosed}.
We return this old field separately rather than treating it as an internal
field of $H_r$, whose internal indices stop at $r-2$.  This proves all five
items.
\end{proof}
